\documentclass[aspectratio=43]{beamer}

% Tema geral
\usetheme{Madrid}
\usecolortheme{default}

% Pacotes
\usepackage[utf8]{inputenc}
\usepackage[T1]{fontenc}
\usepackage[brazil]{babel}
\usepackage{lmodern}

% Cores principais
\definecolor{ufalblue}{RGB}{45, 45, 170}
\definecolor{darkred}{RGB}{170, 0, 0}
\definecolor{darkgreen}{RGB}{0, 110, 20}

% Estilo das barras
\setbeamercolor{palette primary}{bg=ufalblue, fg=white}
\setbeamercolor{palette secondary}{bg=ufalblue, fg=white}
\setbeamercolor{palette tertiary}{bg=ufalblue, fg=white}
\setbeamercolor{frametitle}{bg=ufalblue, fg=white}
\setbeamercolor{title}{fg=white, bg=ufalblue}

% Blocos
\setbeamercolor{block title}{bg=ufalblue, fg=white}
\setbeamercolor{block body}{bg=blue!5, fg=black}

\setbeamercolor{block title alerted}{bg=darkred, fg=white}
\setbeamercolor{block body alerted}{bg=red!8, fg=black}

\setbeamercolor{block title example}{bg=darkgreen, fg=white}
\setbeamercolor{block body example}{bg=green!8, fg=black}

% Rodapé parecido com o do PDF
\setbeamertemplate{footline}{
  \leavevmode%
  \hbox{%
    \begin{beamercolorbox}[wd=.34\paperwidth,ht=2.5ex,dp=1ex,left]{palette primary}
      \hspace{1em}Disciplina de Ciência de Dados (UFAL)
    \end{beamercolorbox}%
    \begin{beamercolorbox}[wd=.33\paperwidth,ht=2.5ex,dp=1ex,center]{palette primary}
      Projeto Final
    \end{beamercolorbox}%
    \begin{beamercolorbox}[wd=.33\paperwidth,ht=2.5ex,dp=1ex,right]{palette primary}
      \insertframedate{}\hspace{1em}\insertframenumber{} / \inserttotalframenumber\hspace{1em}
    \end{beamercolorbox}%
  }
}

% Bolinhas dos itens
\setbeamertemplate{items}[circle]

% Informações da capa
\title{Ciência de Dados: Projeto Final}
\subtitle{Especificações da Disciplina}
\author{Disciplina de Ciência de Dados}
\institute{Universidade Federal de Alagoas}
\date{19 de maio de 2026}

\begin{document}

\begin{frame}
  \titlepage
\end{frame}

\begin{frame}{Sumário}
  \tableofcontents
\end{frame}

\section{Diretrizes e Especificações do Projeto}

\begin{frame}{Objetivos do Projeto Final}

\begin{block}{Proposta Principal}
Aplicar técnicas de Ciência de Dados em uma base de dados escolhida pela equipe,
com foco na geração de conhecimento útil para apoiar uma tomada de decisão real
ou simulada.
\end{block}

O trabalho é avaliado sob três pilares fundamentais:

\begin{enumerate}
  \item \textbf{Conteúdo Técnico:} domínio das etapas de análise, tratamento,
  modelagem e avaliação dos dados.
  \item \textbf{Aplicação Prática:} uso dos resultados para apoiar decisões.
  \item \textbf{Comunicação em Formato de Pitch:} apresentação convincente.
\end{enumerate}

\end{frame}

\begin{frame}{Estrutura do Documento Técnico}

\begin{alertblock}{Padrão de Formatação}
O documento deve seguir o modelo de artigos da SBC e ter limite máximo de 12 páginas.
\end{alertblock}

\begin{enumerate}
  \item \textbf{Aplicação:} proposta, problema, justificativa e objetivos.
  \item \textbf{Base de Dados:} origem, contexto, registros e pré-processamento.
  \item \textbf{Resultados e Discussão:} resultados, limitações e decisão.
\end{enumerate}

\end{frame}

\begin{frame}{Formato do Insight Acionável}

\begin{exampleblock}{Estrutura Recomendada}
Percebemos que \textit{[padrão nos dados]}. Isso sugere que
\textit{[interpretação do problema]}. Por isso, recomendamos
\textit{[ação prática]}.
\end{exampleblock}

\end{frame}

\end{document}